% BEGIN PROFESSOR FEEDBACK PF12: Method - Detection methods
\section{Detection methods}
\label{sec:detection_methods}

The previous two sections produced two descriptions of every address: the
non-motif baseline of Section~\ref{sec:graph_construction} and the motif
fingerprint of Section~\ref{sec:motif_extraction}. This section describes how
those descriptions are turned into a detector and which models are used. The
binary detection for RQ1 and RQ2 is kept apart from the category detection for
RQ3 and RQ4, because the two groups ask different questions and use different
datasets.

\subsection{The three feature views}

RQ1 asks whether the motif fingerprint adds anything beyond ordinary address
behaviour. The cleanest way to answer this is a controlled comparison, in which
only the features change and everything else stays fixed. Three views of each
address are therefore built from the two descriptions above:
\begin{itemize}
  \item the \emph{baseline} view, the 17 non-motif features of
        Table~\ref{tab:results_baseline_features},
  \item the \emph{motifs} view, the log-transformed local motif-role counts, with 32 roles
        in two HuggingFace samples and all 40 roles in the sample that includes
        triangles, and
  \item the \emph{combined} view, the two simply placed side by side.
\end{itemize}
As described in Section~\ref{sec:evaluation_protocol}, RQ1 and RQ2 use three
non-overlapping HuggingFace samples. Two samples use the 32 star and two-node
roles. The remaining sample has the released-neighbourhood closure needed for all
40 roles and evaluates the 32-role view and the eight triangle roles as separate
ablations.
All three views use exactly the same addresses and the same train and test split,
and no address is dropped just because its fingerprint is all zeros. Because the data is
held constant, this comparison evaluates the predictive value added by the
fingerprint representation. The main contrast is 17 baseline features against
those same 17 features plus 32 or 40 local motif-role features. Each role
combines structure, direction, focal-node position, and temporal order.
The experiment therefore does not isolate the contribution of temporal order
alone. That would require a matching representation with the ordering removed.

RQ2 then asks how the time window affects the usefulness of the fingerprints.
For this the
combined view is evaluated at each of the nine windows from
Section~\ref{sec:motif_extraction}. The baseline does not depend on the window, so
it stays the same throughout and serves as a fixed reference line.

\subsection{Classifiers for RQ1 and RQ2}

The detector for the binary questions is a random forest, an ensemble of decision
trees that each vote on whether an address is malicious. A random forest suits
this task because it can pick up interactions between different motif roles
without any manual feature scaling, and because it is a strong, well-understood
default rather than a model tuned hard for the highest possible score. Its
settings are fixed in advance, with 100 trees, a limited depth, balanced class
weights, and a fixed random seed, and they are not adjusted on the test data. The
forest is trained on the development addresses and then scores each address in the
test set with a probability of being malicious. The time window is chosen on the
development data by cross-validation (CV), using the rule described in
Section~\ref{sec:evaluation_protocol}.

As a robustness check, an Extreme Gradient Boosting (XGBoost) model is also fitted, on
the combined view at the chosen window. If it agrees with the random forest, that
strengthens confidence in the result, but the RQ1 decision always rests on the
random forest, which is fixed beforehand. The models use the allocated central
processing unit (CPU) threads. Table~\ref{tab:results_models}
summarises all three models.

\begin{table}[ht]
\centering
\caption{The three models used in the detection experiments.}
\label{tab:results_models}
\small
\begin{tabular}{@{}p{2.8cm} p{3.7cm} p{6.6cm}@{}}
\toprule
\textbf{Task} & \textbf{Model} & \textbf{Fixed settings} \\
\midrule
RQ1 and RQ2 (main) &
Random forest &
100 trees, depth 10, balanced class weights, a fixed seed, allocated CPU threads \\[2pt]
RQ1 (robustness) &
Gradient-boosted trees (XGBoost) &
100 rounds, depth 6, learning rate 0.1, positive-class weight from the training
split, a fixed seed \\[2pt]
RQ3 and RQ4 &
Multinomial logistic regression &
Standardised features, balanced class weights, up to 2{,}000 iterations, a fixed
seed \\
\bottomrule
\end{tabular}
\end{table}

\subsection{Classifier for RQ3 and RQ4}

RQ3 and RQ4 do not separate malicious from benign addresses. They sort criminal
addresses into their behaviour category, so the model has one class per category
rather than two. Both questions are answered from the same model on the same
data, and they are therefore always prepared and run together as one analysis.
Only categories with at least 100 addresses take part, because the categories are
very unevenly filled, as explained in Section~\ref{sec:data_acquisition}, and
categories are never merged into larger groups.

The input here is the per-layer fingerprint of
Section~\ref{sec:motif_extraction}, expressed as local role proportions, which captures
the mix of roles an address occupies in each transaction layer. The model is a
logistic regression in its multi-class form. Its inputs are first put on a common
scale, and this scaling is learned only from the training data so that the test
data cannot influence it. A logistic regression is chosen deliberately, because it
is a simple model with few parameters. This tests whether the motif proportions
alone can separate the categories, without a more complex model that might overfit
the smaller categories.

RQ3 evaluates category separability using motif fingerprints alone. Class
shares and a permutation reference provide context for this separation.
There is no comparison with the conventional behavioural features for category
classification, so RQ3 does not test whether fingerprints outperform those
features.

Two kinds of window are selected on the development data, one for each question. A
single global window is chosen to separate all categories as well as possible at
once, and this is the window RQ3 uses to ask whether the categories have distinct
fingerprints at all. In addition, a separate window is chosen for each category,
the one at which that category is recognised best. Whether these category-specific
windows actually beat the single global window is what RQ4 tests. Both kinds of
window are chosen on the development data only. The final model is refitted on all
development addresses and gives each address in the test set one probability per
category.

\subsection{Handling class imbalance}

The classes are uneven to different degrees across the experiments, only
slightly for the main dataset and strongly for the Real-CATS categories. The
same applies if the optional BLTE run is performed. This is handled in a single
way, by weighting each class in inverse proportion to how often it appears, so
that a rare class counts for as much as a common one during training. The random
forest and the logistic regression use balanced class weights, and the
gradient-boosted model uses the matching positive-class weight from its training
split. No artificial examples are added, and the number of addresses in each
class is left unchanged.
% END PROFESSOR FEEDBACK PF12

